\section{Audit infrastructure for the structural conditions}
\label{sec:audits}

Four operational structural conditions are added to
\citep[Deployment Safety]{lasser2026paper10}'s deployment claim:
(C7.RATE) (upper exposure-rate cap), (REP) (representativeness),
(DISP) (bounded dispersion), and (COAL) (coalition closure);
plus a Lipschitz embedding compatibility precondition
(Assumption~\ref{ass:embedding-companion}) that certifies $\phi$
before (REP) or (DISP) can be evaluated in the embedded
representation. Each admits an operational audit, paralleling the
existing audit structure of \citep[\S 7]{lasser2026paper10}: each
audit produces an attestation on the verification
ledger~\citep{lasser2026exo}, with cadence calibrated to
deployment epoch boundaries. Confidence-bound construction in the
(REP), (DISP), and (COAL) procedures below follows standard
one-sided upper-confidence-bound machinery~\citep{lehmann2005testing}.

\subsection{(C7.RATE) and channel-projection audit}

This audit certifies the two structural inputs to
Theorem~\ref{thm:c7-corollary} other than (C5.HOEFF) and
(C5.MULT) (which are part of \citep[Audit~7]{lasser2026paper10}):
the upper exposure-rate cap $\lammax$, and the verification
protocol's channel-projection structure.

\paragraph{Procedure (rate cap).}
\begin{enumerate}
\item Identify the per-step exposure event sources: ledger commit
  rate, action-execution rate, governance-event rate.
\item Inspect the hard mechanisms enforcing the rate ceiling ---
  ledger commit caps, action-queue depth bounds, governance
  rate-limit middleware, scheduler quotas --- and identify which
  substrate owns the enforcement (substrate-exclusivity from
  \citep{lasser2026exo}).
\item Compute an upper-bound rate $\lammax$ from the certified
  mechanisms (not from observed historical traffic).
\item Document $\lammax$, the enforcing mechanisms, and the
  substrate ownership in the audit attestation.
\end{enumerate}

Empirical traffic windows over $\taumeta$ may be used as
\emph{monitoring} evidence (e.g., to alert on approach to
$\lammax$), but cannot establish the deterministic deployment-class
upper bound by themselves.

\paragraph{Procedure (channel-projection structure).}
\begin{enumerate}
\item Enumerate the projection maps $\pi_j: A \to V_j \cup
  \{*_j\}$ for $j = 1, \ldots, \Kch$ from the verification
  protocol's event-classification policy.
\item Verify that $\pi_j(a) = *_j$ holds for all actions $a
  \notin A_j$, and that $\ell^{(j)}(a) = 0$ on $A \setminus A_j$
  (zero-outside-engagement convention).
\item Document the pairwise overlap sets $A_{j, j'} = A_j \cap
  A_{j'}$ for use by downstream audits.
\end{enumerate}

\paragraph{Cadence.} Initial certification before deployment
activation; re-certification when verification-protocol
classification policy changes, when rate-limit mechanisms
change, or when SPRT detection observes events at rates
approaching $\lammax$.

\paragraph{Failure mode.} Deployments where the exposure event
rate cannot be bounded by certified mechanisms (e.g.,
agent-controlled rate-limit overrides), or where the channel
projection policy admits $\ell^{(j)}(a) \neq 0$ off $A_j$, fall
outside Theorem~\ref{thm:c7-corollary}'s scope. Mitigation: hard
rate limits enforced by substrate-exclusive infrastructure;
verification-protocol classification redesign to respect the
zero-outside-engagement convention.

\subsection{Embedding $\phi$ certification (precondition for REP and DISP)}

(REP) and (DISP) are stated relative to the embedded deviation
field $\Delta_c = \phi(U_c) - \phi(W)$ in the Hilbert space
$\mathcal{V}$ (Assumption~\ref{ass:embedding-companion}). The
embedding $\phi$ must be certified before (REP) or (DISP) can
be audited.

\paragraph{Procedure.}
\begin{enumerate}
\item Specify the embedding $\phi: \mathcal{V}_\mathrm{utility}
  \to \mathcal{V}$ in deployment documentation.
\item Verify $\phi$ is affine isometric on the operationally
  active subspace $P^\mathrm{act}$: $\|\phi(x) - \phi(y)\| =
  \|x - y\|_{\mathcal{V}_\mathrm{utility}}$ for $x, y \in
  P^\mathrm{act}$, with linearity up to a fixed translation.
\item If $\phi$ is only bi-Lipschitz (not isometric), document
  the bi-Lipschitz constants $(\underline{L}, \bar{L})$ and apply
  them as multipliers to the (REP) / (DISP) bounds downstream.
\end{enumerate}

\paragraph{Failure mode.} If $\phi$ is not bi-Lipschitz on
$P^\mathrm{act}$, the algebraic kernel of
Theorem~\ref{thm:concgap-scoped} cannot be transferred to the
deployment claim's $\|\Pproxy - \Ttruth\|$ form via
Lipschitz transfer~\citep[Theorem~2]{lasser2026micro}. Mitigation:
choose a different embedding, or restrict the deployment's scope
to a subspace where bi-Lipschitz $\phi$ exists.

\subsection{(REP) audit (representativeness)}

\paragraph{Procedure.}
\begin{enumerate}
\item Define the base population measure $\mu$: the auditor's
  fair-or-natural distribution over admissible counterparties
  (typically uniform over S1-admissible
  participants~\citep{lasser2026paper8a}, or weighted by some
  natural allocation).
\item In the embedded representation under the certified $\phi$,
  produce a one-sided upper confidence bound on $\|\bar\Delta\|$
  where $\bar\Delta := \mathbb{E}_\mu[\phi(U_c) - \phi(W)]$. The
  bound is computed by sampling counterparties under $\mu$ (or
  reweighting observed-attestation samples to $\mu$ with stated
  coverage assumptions) from utility-disclosure attestations on
  the verification ledger; conservative bounding is required
  because (REP) is a hypothesis the deployment claim conditions
  on, not a quantity the deployment claim merely monitors.
\item Certify $\|\bar\Delta\| \leq \rho_\mathrm{rep}$ at the
  audit's chosen confidence level for a deployment-class threshold
  $\rho_\mathrm{rep}$.
\item Document $\rho_\mathrm{rep}$, the confidence level, and the
  sampling/reweighting design in the audit attestation.
\end{enumerate}

\paragraph{Cadence.} Initial calibration before deployment;
re-calibration on counterparty-population changes (additions,
removals, or substantial weight-redistribution).

\paragraph{Failure mode.} The counterparty population's mean
utility deviates systematically from the welfare-relevant truth
$W$. Mitigation: counterparty-selection-process diversification,
category-coverage audit~\citep[Audit
3]{lasser2026paper10}, or restriction of the deployment's scope
to a counterparty subset where (REP) holds.

\subsection{(DISP) audit (bounded dispersion)}

\paragraph{Procedure.}
\begin{enumerate}
\item In the embedded representation under the certified $\phi$,
  produce a one-sided upper confidence bound on the
  $\mu$-population variance $\int \|\Delta_c -
  \bar\Delta\|^2 \,d\mu(c)$, where per-counterparty deviations
  $\Delta_c = \phi(U_c) - \phi(W)$ are read from
  utility-disclosure attestations. As with (REP), the variance is
  taken with respect to $\mu$, not the trade-flow weighting $w$:
  audit sampling under $\mu$ (or reweighting to $\mu$ with stated
  coverage) is required.
\item Certify the variance $\leq \sigma^2$ at the audit's chosen
  confidence level for a deployment-class constant $\sigma$.
\item Document $\sigma$, the confidence level, and the
  sampling/reweighting design in the audit attestation.
\end{enumerate}

\paragraph{Cadence.} Continuous attestation under the calibrated
$\mu$-sampling design over a rolling calibration window;
confidence-bound violations trigger audit re-calibration. Observed
trade-flow ($w$) statistics may serve as monitoring signal but do
not by themselves certify the $\mu$-variance hypothesis.

\paragraph{Failure mode.} High dispersion across counterparties
(some counterparties' utilities deviate far from $W$, even when
the population mean is centered). The bound in
Theorem~\ref{thm:concgap-scoped} weakens but does not collapse;
the deployment claim accommodates large $\sigma$ at the cost of
larger Goodhart-slack ceiling.

\subsection{(COAL) audit (coalition closure)}

\paragraph{Procedure.}
\begin{enumerate}
\item Run \emph{coalition-detection analysis} on the verification
  ledger: identify counterparties with correlated trade flows
  above a deployment-class threshold (using
  repeated-beneficiary linkage, governance-vote alignment,
  attestation-source clustering).
\item Partition $\mathcal{C}$ into equivalence classes
  (coalitions) by the detected correlation structure.
\item Re-compute $\HHI$ on the post-partition (coalition-level)
  weights.
\item \textbf{Certify clause~(ii) of (COAL).} Document the
  deployment-class bound $N_{\max}$ on the post-partition
  counterparty cardinality, together with the onboarding /
  governance policy that prevents $N$ from scaling with $|P|$.
  Document the choice of base measure $\mu$ as the uniform
  distribution on the post-partition counterparties (the
  structural prerequisite for the HHI-to-$\chi^2$ translation
  $\chi^2(w \,\|\, \mu) = N \HHI(w) - 1$). Deployments that
  prefer to commit directly to a $\chi^2$ ceiling certify
  $\chi^2(w \,\|\, \mu) \leq \Xi$ for a deployment-class constant
  $\Xi$ instead, with $\mu$ documented as whatever fair-or-natural
  base measure the audit selects.
\item Produce a one-sided upper confidence bound on the
  latent-coalition residual $\eta_\mathrm{latent}$ from the
  audit's detection-power calibration: bound $\eta_\mathrm{latent}
  \leq p_\mathrm{miss} \cdot \|\Delta\|_{\max}$ where
  $p_\mathrm{miss}$ is the threshold-miss probability for the
  worst-case undetected coalition size and $\|\Delta\|_{\max}$ is
  the worst-case per-counterparty deviation magnitude in the
  deployment's threat model.
\item Document the coalition partition, $N_{\max}$ (or $\Xi$),
  and the $\eta_\mathrm{latent}$ confidence bound (with confidence
  level) in the audit attestation.
\end{enumerate}

\paragraph{Cadence.} Continuous coalition detection on the
verification ledger; audit-cadence alerts on detected
coalition-formation events; re-partition triggered when new
coalitions form.

\paragraph{Failure mode.} Many small counterparties acting in
coordinated fashion (a hidden coalition) drive the effective
$\HHI$ above the audited threshold. The deployment claim's
bound on Goodhart slack picks up the $\eta_\mathrm{latent}$
residual; deployments where $\eta_\mathrm{latent}$ cannot be
bounded operationally are outside
Theorem~\ref{thm:concgap-scoped}'s scope.

\subsection{Integration with existing audit infrastructure}

Several of these new audits overlap with existing audits
in~\citep[\S 7]{lasser2026paper10}:
\begin{itemize}
\item (C7.RATE) and channel-projection audit: integrates
  primarily with Audit~4 (bounded co-evolution
  calibration, dual-mode), which is the deployment-tooling home
  for the structural inputs to
  Theorem~\ref{thm:c7-corollary}; Audit~7 (SPRT-applicability +
  gap-growth + clock comparability) supplies the (C5.HOEFF) clip
  radius and related SPRT-side calibration inputs but not the
  (C7.RATE) upper bound.
\item (REP) audit: integrates with Audit~3
  (cooperative-vs-redundancy audit), since both rely on
  category-coverage analysis.
\item (DISP) audit: new, but uses utility-disclosure attestations
  that the existing infrastructure already produces.
\item (COAL) audit: integrates with $I_9$
  (substrate-exclusivity observability) and $I_{10}$
  (coverage/materiality routing) of \citep[Deployment Safety]{lasser2026paper10},
  since coalition detection uses ledger-linkage analysis those
  invariants already specify.
\end{itemize}

The audit hooks above are the minimum needed for
Theorem~\ref{thm:c7-corollary} and Theorem~\ref{thm:concgap-scoped}
to apply; consolidated audit-tooling combining these with
\citep[\S 7]{lasser2026paper10}'s existing audits is the natural
target for further operational-tooling work.
